\section{A general description of the environment} \label{section:general_description}
Every year, the firm must estimate its technical provisions given the information available to it at that point in time. We define technical provisions as the value of funds required to transfer those liabilities (`transfer value').

\subsection{Information set} \label{section:information_set}
We denote the information set as \begin{math} I_t \subset I \end{math} as being the set of all available information, or the `global' information set, at time \begin{math} t \end{math}, where \begin{math} I \end{math} is the set of all possible information sets (the `information universe'). It is then clear that \begin{math} I_t \subset I_{t+j} \forall j \geq 0 \end{math} (i.e. any information known in a given year is also known in all subsequent years). We also denote the set \begin{math} i_t \subset I_t \end{math} as the set of information \textit{known to the firm}, or the 'local' information set, at time \(t\). 

\subsection{The firm's liability portfolio} \label{section:liability_portfolio}
We denote the portfolio set \begin{math} P_t \subset P \end{math} as the set describing the portfolio characteristics of the firm, or the portfolio-specific information, where, as above, \begin{math} P \end{math} is the set of all possible portfolios (the `portfolio universe').

\subsection{Information choice function} \label{section:information_choice_function}
Given this, we denote the set function \begin{math} g_t:i_t \rightarrow i_t \end{math} as the 'information choice function' used within the firm which maps the local information set onto a subset of that. This subset is the information which is used by the firm as an input into its valuation model (which can include expert judgement - see Section \ref{section:valuation}). In other words, this function represents the process used within the firm to select the information used to help formulate a value for the firm's liabilities. Keeping notation consistent, we define \(g\) as the set of all possible information choice functions, so that \begin{math} g_t \subset g \end{math}.

\subsection{Valuation} \label{section:valuation}
We can now define the set function \begin{math} f_t:P \times I \rightarrow \mathbb{R} \end{math} as the function at time \begin{math} t \end{math} which gives the transfer value of the liabilities for portfolio \begin{math} P_t \end{math} given the information set \begin{math} I_t \end{math}.

We distinguish between `true' VaR and `estimated' VaR. We make this distinction since we can never know for certain the true probability law underlying a random process and therefore a firm's VaR can only be an estimate of the true VaR for that risk and as such is a random variable itself. Under this framework, we can then make explicit statements about model risk by comparing the true VaR to a firm's estimated VaR.
%
\subsection{Specifying `true' VaR}
In order to specify the true one-year VaR, we first need to define the `true' transfer value of a firm's technical provisions in one year's time, given the global and the portfolio-specific information at that moment in time; consistent with the notation above, this is given as \begin{math} f_{t+1} \left( P_{t+1}, I_{t+1} \right) \end{math}. We also define the function \begin{math} c:P \times P \rightarrow \mathbb{R} \end{math} as the one-year cashflow of the firm. For notational brevity, we define \begin{math} c_{t} \in \mathbb{R} \end{math} as \begin{math} c_{t} = c \left( P_{t}, P_{t+1} \right)\end{math} and the random variable \begin{math} \tilde{c}_{t} = c \left( P_{t}, \tilde{P}_{t+1} \right)\end{math}

The problem we are then faced with is to find a function \begin{math} g:P \times I \times \left[ 0,1 \right] \rightarrow \mathbb R \end{math} which, at any given time \begin{math} t \end{math}, produces the risk capital for a one-year VaR for a portfolio \begin{math} P_t \subset{P} \end{math} and information set \begin{math} I_t \subset{I} \end{math} at the given confidence level \begin{math} k \in{\left[ 0, 1 \right]} \end{math} (e.g. \begin{math} k = 0.995 \end{math} as in the ICA/SCR). This implies that the probability \begin{math} K \end{math} of remaining solvent over the run-off of the portfolio (i.e. the probability of paying all liabilities in full over the lifetime of the portfolio) is given as \begin{math} K = \left( 1 - k \right) ^ T \end{math}, where \begin{math} T \end{math} is the term of the run-off\footnote{Note that these conditions combined are stronger than simply the condition that the probability of non-ruin over the portfolio run-off is K, which can be achieved by setting the probability of non-ruin over some years to be lower than \begin{math} k \end{math} and over other years to be higher than \begin{math} k \end{math}}.

\subsection{Solution to the problem}
We now define the function \begin{math} g:P \times I \times \left[ 0,1 \right] \rightarrow \mathbb R \end{math} as the function which, at time \begin{math} t \end{math}, produces the firm's risk capital for portfolio \begin{math} P_t \subset{P} \end{math} and information set \begin{math} I_t \subset{I} \end{math} at the given confidence level \begin{math} k \in{\left[ 0, 1 \right]} \end{math}.
\begin{definition}[\begin{math} A \left( y \right) \end{math}]
The multivalued function \begin{math} A:\mathbb{R} \rightarrow \mathbb{R} \times I \times F \times F \end{math} produces a set of values such that:
\begin{equation}
\left[ \tilde{c}_t - \mathrm{E} \left[ c_t \right] \right]
+ \left[ f \left( \tilde{I}_{t+1} \right) - \mathrm{E} \left[ f \left( I_{t+1} \mid I_t \right) \right] \right] \leq k
\end{equation}
\end{definition}
\begin{equation}
\Pr \left\{ \left( \left[ \tilde{c}_t - \mathrm{E} \left[ c_t \right] \right]
+ \left[ f \left( \tilde{I}_{t+1} \right) - \mathrm{E} \left[ f \left( I_{t+1} \mid I_t \right) \right] \right] \right) \cdot \left( 1 + r_t \right) ^ {-1} \geq g \left( P_t, I_t, k \right) \right\}
= k
\end{equation}
where \begin{math} r_t \in{\mathbb{R ^ +}} \end{math} is the one-year spot risk-free rate of interest at time \begin{math} t \end{math}. The SCR at time \begin{math} t \end{math} is therefore given as \begin{math} g \left( P_t, I_t, 0.005 \right) \end{math}
%
\begin{equation}
\left( \hat{c}^*, \hat{I}^*_{t+1}, \hat{f}^*_{t+1}, \hat{F}^*_{t+1} \right)
\in \argmin_{c, I_{t+1}, f_{t+1}, F_{t+1}}
\left( \left[ \tilde{c}_t - \mathrm{E} \left[ c_t \right] \right]
+ \left[ f \left( \tilde{I}_{t+1} \right) - \mathrm{E} \left[ f \left( I_{t+1} \mid I_t \right) \right] \right] \right) \cdot \left( 1 + r_t \right) ^ {-1} \geq g \left( P_t, I_t, k \right)
\end{equation}